% =============================================================================
% Semana 12 — Trabalho 6 (Apresentações Finais) + Revisão para Prova II
% Aula de 3 horas (19h–22h) | 20/05/2026
% =============================================================================

\chapter{Semana 12 (20/05/2026) — Apresentações Finais e Revisão para Prova II}

% ---------------------------------------------------------------------------
\section{Objetivo da semana}
% ---------------------------------------------------------------------------

Pessoal, chegamos à última aula da disciplina. Hoje é o dia das
\textbf{apresentações finais do Trabalho 6} — o projeto mais completo do
semestre, integrando AED + Storytelling + Tableau. Depois fechamos com
uma revisão integrada da Unidade II para a Prova Final.

Esta aula é a síntese de tudo: do dado bruto até a decisão comunicada.

\begin{BoardBox}
\textbf{Roteiro da aula (19h–22h)}
\begin{enumerate}
  \item 19h–20h30 — \textbf{Apresentações:} Trabalho 6 (pitch 3 min + demo) — todos os alunos
  \item 20h30–21h — \textbf{PEER REVIEW:} feedback estruturado coletivo
  \item 21h–21h50 — \textbf{REVISÃO:} Unidade II — simulado final
  \item 21h50–22h — \textbf{Fechamento da disciplina:} reflexão e próximos passos
\end{enumerate}
\end{BoardBox}

% =============================================================================
\section{Parte I — Apresentações Finais T6 (19h–20h30)}
% =============================================================================

\begin{BoardBox}
\textbf{Rubrica de avaliação T6 — Dashboard Final (para o quadro)}

\begin{tabular}{|p{4cm}|r|p{8cm}|}
\hline
\textbf{Critério} & \textbf{Peso} & \textbf{Descrição} \\
\hline
Big Idea + Pergunta de Decisão & 15\% & Clara, com ponto de vista e tensão; explicitada no pitch \\
Estrutura narrativa & 15\% & Contexto → diagnóstico → recomendação → call to action \\
Qualidade dos visuais & 20\% & Títulos-mensagem; design sem chartjunk; hierarquia visual \\
Interatividade & 15\% & Parâmetro + filter action funcionando; comportamento intuitivo \\
KPIs com targets & 15\% & Metas explícitas; gap calculado; indicador de status \\
Pitch 3 minutos & 10\% & Dentro do tempo; dados corretos; recomendação concreta \\
Dashboard publicado & 10\% & Link Tableau Public funcional antes da aula \\
\hline
\textbf{Total} & \textbf{100\%} & \\
\hline
\end{tabular}
\end{BoardBox}

\begin{SolvedBox}
\textbf{Ordem das apresentações:} sorteio na abertura da aula.\\
\textbf{Timing:} 3 min de pitch + 2 min de perguntas da turma + 1 min de feedback do professor.\\
\textbf{Link do Tableau Public:} enviado no chat antes de começar — professor testa ao vivo.

\textbf{O que acontece se o link não funcionar?} A apresentação ocorre com screenshots
e o aluno perde os 10\% de bônus de publicação. O trabalho não é invalidado.

\textbf{Durante as apresentações,} cada integrante da turma preenche um
formulário de observação (3 linhas por apresentação):
\begin{enumerate}
  \item Big Idea em 1 frase (copiada/adaptada do que ouviu)
  \item 1 decisão que tomaria com base no dashboard
  \item 1 sugestão de melhoria
\end{enumerate}

Isso exercita \emph{escuta ativa} ao mesmo tempo que gera feedback de valor para o apresentador.
\end{SolvedBox}

% =============================================================================
\section{Parte II — Peer Review Coletivo (20h30–21h)}
% =============================================================================

\begin{BoardBox}
\textbf{Sessão de feedback estruturado — dinâmica de galeria}

\begin{enumerate}
  \item Todos os dashboards ficam abertos em abas no projetor (ou em laptops compartilhados).
  \item Cada aluno anota em post-it (virtual ou físico) 1 ponto forte e 1 sugestão
  para \textbf{pelo menos 3} trabalhos que não o seu.
  \item Professor facilita um round de highlights: ``O que mais apareceu como ponto forte?
  O que mais apareceu como melhoria?''
  \item Identificar 2–3 padrões de excelência e 2–3 padrões de problema recorrente
  para encerrar a atividade com aprendizado coletivo.
\end{enumerate}
\end{BoardBox}

% =============================================================================
\section{Parte III — Revisão Unidade II (21h–21h50)}
% =============================================================================

\begin{SolvedBox}
\textbf{Mapa mental da Unidade II — para o quadro}

\textbf{Semana 08:} AED vs Explanatória; Big Idea framework (ponto de vista + tensão + stake);
pergunta de decisão; storyboard 4–6 telas.\\
\textbf{Semana 09:} Pré-atenção (cor, tamanho, posição); Teste do Relance;
título-mensagem; Chartjunk (Tufte) + data-ink ratio; hierarquia visual; consistência.\\
\textbf{Semana 10:} Tableau — filtros (hierarquia e tipos); parâmetros (criação e uso
em campos calculados); ações (filter, highlight, URL, navigate).\\
\textbf{Semana 11:} Dashboards decisórios (contexto→diagnóstico→recomendação);
KPIs com target + gap + tendência; bullet chart; consistência de definições;
pitch de 3 minutos (situação → complicação → resolução → próximo passo).\\
\textbf{Semana 12:} Feedback, refinamento, publicação e apresentação executiva.
\end{SolvedBox}

\begin{BoardBox}
\textbf{Simulado Prova II — 8 questões (30 min para responder + 15 min correção)}

\textbf{Q1.} Qual a diferença entre análise exploratória e análise explanatória?
Dê um exemplo de quando cada uma é a abordagem correta.

\textbf{Q2.} Escreva uma Big Idea completa para o seguinte cenário:
``O tempo médio de resposta do suporte online passou de 4h para 11h no último trimestre,
enquanto o NPS caiu de 72 para 58.''

\textbf{Q3.} Um dashboard tem 12 KPIs na tela inicial, todos sem metas,
todos em verde. Identifique 3 problemas e proponha como corrigi-los.

\textbf{Q4.} Explique o que é o Teste do Relance. Um gráfico falha no teste.
Liste 3 causas possíveis e como resolver cada uma.

\textbf{Q5.} Qual a diferença entre um Quick Filter e um Parâmetro no Tableau?
Dê um exemplo em que o parâmetro é necessário (algo que o filtro não faz).

\textbf{Q6.} Você está apresentando para o Conselho Diretivo da empresa e descobriu
que o canal de vendas diretas gera 3× mais receita por cliente do que o e-commerce,
mas representa apenas 12\% dos clientes. Monte o storyboard de 4 telas:
título de cada tela + tipo de visual.

\textbf{Q7.} Um analista mostra um gráfico de pizzahcom 8 fatias em 8 cores diferentes.
O título é ``Distribuição de Produtos 2025''. Quais são os problemas? Redesenhe
descrevendo em texto como ficaria o visual correto.

\textbf{Q8.} Qual a estrutura do pitch de 3 minutos? Escreva os 4 blocos com
exemplos para o cenário: empresa de logística com 40\% dos pedidos atrasando
por falta de otimização de rota.
\end{BoardBox}

\begin{SolvedBox}
\textbf{Gabaritos do Simulado (professor):}

\textbf{Q1.} Exploratória: analista descobrindo padrões no dado para si mesmo —
centenas de visuais gerados, a maioria descartada, linguagem neutra. Ex.: investigar um
novo dataset antes de saber o que está procurando. Explanatória: analista
comunicando uma descoberta para provocar uma decisão específica — apenas os
visuais que suportam a mensagem, linguagem de impacto. Ex.: apresentação para
o CEO sobre oportunidade de redução de custos.

\textbf{Q2.} \textit{Exemplo:} ``Se não corrigirmos o processo de atendimento online
nos próximos 30 dias, o NPS continuará caindo — e com ele a renovação dos contratos
anuais que vencem no Q2.'' (ponto de vista + tensão + stake)

\textbf{Q3.} Problemas: (1) 12 KPIs excede capacidade atencional — onde olhar primeiro?
(2) Sem meta, o número não tem significado — R\$10.000 é bom ou ruim? (3) Tudo verde
cria falsa sensação de controle. Correções: (1) máximo 5 KPIs na tela inicial, o resto em
drill-down; (2) cada KPI com target e gap calculado; (3) vermelho/amarelo/verde baseado
em desvio real da meta.

\textbf{Q4.} Teste do Relance: mostrar o visual por 3 segundos; perguntar o que a pessoa
entendeu. Causas de falha: (a) título rótulo — correção: título mensagem; (b) excesso de
elementos — correção: remover tudo que não é dado; (c) cor sem propósito — correção:
1 cor de destaque para o elemento que carrega a mensagem.

\textbf{Q5.} Quick Filter: inclui/exclui linhas do dataset (filtra dados). Parâmetro:
muda o \emph{cálculo} sem filtrar linhas. Ex. que filtro não faz: mudar a métrica
visualizada de Receita para Margem — o filtro não troca o campo; o parâmetro permite
que o usuário selecione qual campo calcular.

\textbf{Q6.}
\begin{enumerate}
  \item ``Direto é 3× mais lucrativo por cliente'' — KPI comparativo: barra Direto vs E-commerce com rótulo R\$/cliente
  \item ``Mas representa apenas 12\% dos clientes'' — donut ou barra de composição do mix de canais
  \item ``O que aconteceria se aumentássemos para 25\%?'' — parâmetro de simulação de mix + KPI de impacto estimado
  \item ``Recomendamos redirecionar a força de vendas para o canal direto até setembro'' — texto + call to action
\end{enumerate}

\textbf{Q7.} Problemas: (a) 8 cores — impossível distinguir mais de 6; (b) pizza com muitas fatias —
olho não compara ângulos com precisão; (c) título rótulo sem mensagem. Redesenho:
barras horizontais ordenadas (maior$\to$menor); 1 cor de destaque para o produto líder;
título: ``Produto C concentra 38\% das vendas''; rótulos diretamente nas barras.

\textbf{Q8.}
\begin{enumerate}
  \item \textbf{Situação:} ``40\% dos pedidos entregamos com atraso — custo de R\$2,5M em multas e créditos no último ano.''
  \item \textbf{Complicação:} ``A análise mostra que 78\% desses atrasos ocorrem em 3 rotas específicas do interior de SP, que usamos há 5 anos sem revisão.''
  \item \textbf{Resolução:} ``Com otimização de rota baseada em dados de trânsito em tempo real, estimamos reduzir os atrasos em 60\% — economizando R\$1,5M ao ano. O investimento no sistema é R\$180k.''
  \item \textbf{Próximo passo:} ``Aprovação de R\$180k até próxima sexta para iniciar o piloto nas 3 rotas críticas no mês que vem.''
\end{enumerate}
\end{SolvedBox}

% =============================================================================
\section{Parte IV — Fechamento da Disciplina (21h50–22h)}
% =============================================================================

\begin{BoardBox}
\textbf{Reflexão final — 3 perguntas para a turma (para o quadro)}

\begin{enumerate}
  \item \textit{``Qual foi o maior insight que você teve ao longo do semestre — não sobre
  os dados, mas sobre como você \textbf{pensa} sobre dados?''}
  \item \textit{``Se você entrasse amanhã numa equipe de dados, qual das ferramentas
  ou metodologias que aprendemos você usaria primeiro?''}
  \item \textit{``O que você ainda quer aprender sobre análise de dados que não cobrimos
  aqui?''} (registrar respostas — direciona planejamento futuro da disciplina)
\end{enumerate}
\end{BoardBox}

\begin{NoteBox}
\textbf{Para a Prova II (21/05 ou conforme calendário oficial):}
\begin{itemize}
  \item Conteúdo: Unidade II completa (storytelling, design, Tableau, dashboards, pitch)
  \item Formato: questões abertas e análise de caso
  \item Tableau: pode aparecer screenshot de dashboard para diagnóstico; sem necessidade
  de acessar o software durante a prova
  \item Big Idea e pitch: podem ser solicitados por escrito
  \item Revisão: simulado desta semana é o melhor guia de estudo
\end{itemize}
\end{NoteBox}

\begin{SolvedBox}
\textbf{Checklist do Trabalho 6 — entrega final até 20/05 às 23h59:}
\begin{itemize}
  \item[$\square$] Link Tableau Public funcional (testado em modo privado $\to$ público)
  \item[$\square$] Mínimo 3 sheets no dashboard; parâmetro + filter action funcionando
  \item[$\square$] Todos os títulos são mensagens (Teste do Relance aprovado)
  \item[$\square$] KPIs com meta, gap e indicador de status
  \item[$\square$] Big Idea visível no dashboard
  \item[$\square$] Pitch de 3 min preparado e cronometrado
  \item[$\square$] Relatório: 1 página com Big Idea + dados usados + decisão recomendada
\end{itemize}
\end{SolvedBox}

% =============================================================================
\section{Revisão Semanal — Questão tipo Prova II (Q5)}
% =============================================================================

\begin{NoteBox}
Última revisão antes da \textbf{Prova II}. Tema: pitch de 3 minutos
e apresentação de resultados.
\end{NoteBox}

\begin{SolvedBox}
\textbf{Cenário (estilo Q5 da Prova II):}

Um aluno de Sistemas de Informação vai apresentar sua análise sobre
evasão em cursos de tecnologia no Agreste de Pernambuco. Ele tem
\textbf{3 minutos} e estruturou o pitch assim:

\textbf{Minuto 1 — Problema (gancho):}
``A cada semestre, 28\% dos alunos de TI no Agreste trancam a matrícula.
Isso custa R\$ 4,2 milhões em bolsas e infraestrutura ociosa. Nosso
modelo identificou os 3 fatores mais preditivos de evasão.''

\textbf{Minuto 2 — Insights (dados):}
\begin{itemize}
  \item Visualização 1: heatmap de evasão por período e turno
  \item Visualização 2: correlação entre frequência no 1º período e
  permanência
  \item Insight-chave: ``Alunos com frequência $<$ 70\% no 1º período
  têm 5$\times$ mais chance de evadir.''
\end{itemize}

\textbf{Minuto 3 — Recomendação (call-to-action):}
``Propomos tutoria intensiva no 1º período para alunos com frequência
abaixo de 70\%. Projeção: redução de 20\% na evasão, beneficiando
850 alunos por ano. Próximo passo: piloto em 3 campi no semestre 2026.2.''

\textbf{Questão:} identifique \textbf{4 elementos} que tornam esse
pitch eficaz.

\textbf{Resolução:}
\begin{enumerate}
  \item \textbf{Gancho com número impactante:} ``28\% trancam''
  + custo em reais captura atenção imediata.
  \item \textbf{Apenas 2 visualizações:} não sobrecarrega — mostra
  só o essencial para sustentar o argumento.
  \item \textbf{Insight traduzido em linguagem acionável:} não diz
  ``$r = 0{,}65$'', diz ``5$\times$ mais chance'' — decisão, não
  estatística.
  \item \textbf{Call-to-action com métricas e próximo passo:}
  ação concreta (tutoria), impacto quantificado (850 alunos/ano),
  e prazo (piloto 2026.2) — o público sabe exatamente o que fazer.
\end{enumerate}
\end{SolvedBox}

% ---------------------------------------------------------------------------
\section{Materiais Preparados}
% ---------------------------------------------------------------------------
\begin{itemize}
  \item Formulário de observação de pitch (impresso — 1 por aluno)
  \item Rubrica T6 projetada durante as apresentações
  \item Simulado impresso (1 por aluno)
  \item Cronômetro visível para os pitches
\end{itemize}
